%! Characteristics of Functions: Increasing, Decreasing, Extrema, and Concavity — Unit 2, Lesson 2.5 — AP Practice
% EXTRA CREDIT — optional, exactly TWO pages, placed between the notes and the graded homework.
% Page 1 is Section I, five multiple-choice items with four options (A)–(D); page 2 is Section
% II, one multi-part free-response question on a pre-drawn graph. Its contexts (a lake's fish,
% a harbor's tide, an unlabeled h) are used nowhere else in the lesson.
% Distractors are the lesson's errors on purpose: MC 1 swaps the value and its location (C)
% or reads a relative max as a span of time (D); MC 2 equates concave down with decreasing (B,
% D) or increasing with concave up (A); MC 3 reports the interval in y-values (B); MC 4
% calls a turning point an inflection point (A, B); MC 5 gives the x-location (B, C) or the
% relative max (A) as the absolute maximum value.
%
% PAGE LOCKSTEP with the other half of the pair: the key marks the correct option in its
% LABEL (no text added to the line) and prose answers are \writespace{H}{…}, fixed height both
% ways. The explicit \newpage fixes the page break in both files.
\documentclass[10pt]{article}
\usepackage{precalculus-article}
\usepackage{precalculus-boxes}
\newcommand{\TallMath}[1]{$\displaystyle #1\rule[-1.4em]{0pt}{3.2em}$}

\begin{document}

\pageheader{Unit 2, Lesson 2.5}{AP Practice}

\vspace{0.12in}

\boxguard[8]
\begin{remindbox}
\small
\textbf{Extra credit --- optional.} These two pages are written in the style of the AP
Precalculus exam. Your graded homework is the last section of this packet; do it first. Every
answer choice below is a mistake someone really makes, so read all four before you choose.
\end{remindbox}

\vspace{0.12in}

\begin{headlinebox}{goldbg}
\small
\textbf{Section I --- Multiple Choice.} Circle the letter of the single best answer. No calculator.
\end{headlinebox}

\vspace{0.06in}

\begin{enumerate}[leftmargin=1.9em, itemsep=12pt]

\item The function $F$ gives the number of fish in a lake $t$ years after 2020. The graph of
      $F$ has a relative maximum at the point $(6, 1200)$. Which of the following is the best
      interpretation?
      \begin{enumerate}[leftmargin=2.6em, itemsep=2pt]
        \item[(A)] The greatest number of fish the lake ever holds is 6.
        \item[(B)] Six years after 2020 the population peaks at 1200 fish, more than in the years
                   just before and just after.
        \item[(C)] 1200 years after 2020, the lake holds 6 fish.
        \item[(D)] The fish population stays at 1200 for 6 years.
      \end{enumerate}

\item The function $g$ is increasing and concave down on the interval $(2, 5)$. Which of the
      following describes the values of $g$ on that interval?
      \begin{enumerate}[leftmargin=2.6em, itemsep=2pt]
        \item[(A)] They increase, faster and faster.
        \item[(B)] They decrease, more and more slowly.
        \item[(C)] They increase, more and more slowly.
        \item[(D)] They decrease, faster and faster.
      \end{enumerate}

\item[] \begin{minipage}[t]{0.5\linewidth}
      \textbf{Questions 3--5 refer to the graph of $h$}, shown for $-2 \le x \le 6$ with both
      endpoints included.

      \vspace{8pt}
      \begin{enumerate}[label=\arabic*., leftmargin=1.9em, itemsep=10pt, start=3]
      \item On which interval is $h$ decreasing?

            \smallskip
            \begin{tabular}{@{}p{3.0cm}p{3.0cm}@{}}
            \textbf{(A)} $(0, 4)$ & \textbf{(B)} $(-1, 3)$ \\[2pt]
            \textbf{(C)} $(4, 6)$ & \textbf{(D)} $(-2, 0)$
            \end{tabular}

      \item Which point on the graph of $h$ is a point of inflection?

            \smallskip
            \begin{tabular}{@{}p{3.0cm}p{3.0cm}@{}}
            \textbf{(A)} $(0, 3)$ & \textbf{(B)} $(4, -1)$ \\[2pt]
            \textbf{(C)} $(2, 1)$ & \textbf{(D)} $(6, 5)$
            \end{tabular}

      \item What is the absolute maximum value of $h$?

            \smallskip
            \begin{tabular}{@{}p{3.0cm}p{3.0cm}@{}}
            \textbf{(A)} $3$ & \textbf{(B)} $6$ \\[2pt]
            \textbf{(C)} $0$ & \textbf{(D)} $5$
            \end{tabular}
      \end{enumerate}
      \end{minipage}\hfill
      \begin{minipage}[t]{0.44\linewidth}
      \vspace{-4pt}
      \centering
      \begin{tikzpicture}[x=0.6cm, y=0.5cm, baseline=(current bounding box.north)]
        \draw[linegray, very thin, step=1] (-3,-2) grid (7,6);
        \draw[->, thick] (-3,0) -- (7.5,0) node[right] {\scriptsize $x$};
        \draw[->, thick] (0,-2) -- (0,6.5) node[above] {\scriptsize $h(x)$};
        \foreach \x in {-2,-1,1,2,3,4,5,6} \node[below, font=\tiny] at (\x,0) {\x};
        \foreach \y in {-1,1,2,3,4,5} \node[left, font=\tiny] at (0,\y) {\y};
        \draw[very thick, plum] (-2,1) .. controls (-1.5,2.2) and (-0.8,3) .. (0,3)
           .. controls (2,3) and (2,-1) .. (4,-1) .. controls (4.8,-1) and (5.4,1.5) .. (6,5);
        \fill[plum] (-2,1) circle (2.5pt);
        \fill[plum] (6,5) circle (2.5pt);
      \end{tikzpicture}
      \end{minipage}

\end{enumerate}

\newpage

\begin{headlinebox}{goldbg}
\small
\textbf{Section II --- Free Response.} Show your reasoning. Answers about the harbor must be
written \emph{in context}, with units --- that is what earns credit on the AP exam.
\end{headlinebox}

\vspace{0.08in}

\begin{enumerate}[leftmargin=1.9em, itemsep=10pt, start=6]

\item The graph shows $W(t)$, the depth of the water at a harbor's dock, in feet, $t$ hours
      after midnight, for $0 \le t \le 12$.

      \begin{center}
      \begin{tikzpicture}[x=0.85cm, y=0.42cm, >=stealth]
        \draw[linegray, very thin, step=1] (0,0) grid (12,9);
        \draw[->, thick] (0,0) -- (12.6,0) node[below right] {\small $t$};
        \draw[->, thick] (0,0) -- (0,9.6) node[above] {\small $W(t)$};
        \foreach \x in {1,...,12} \node[below, font=\scriptsize] at (\x,0) {\x};
        \foreach \y in {2,4,6,8} \node[left, font=\scriptsize] at (0,\y) {\y};
        \draw[very thick, plum, domain=0:12, samples=97, smooth]
             plot (\x, {5 + 3*sin(30*\x)});
        \fill[plum] (0,5) circle (2.5pt);
        \fill[plum] (12,5) circle (2.5pt);
      \end{tikzpicture}
      \end{center}

      \textbf{(a)} On which intervals is $W$ increasing? On which is it decreasing?
      \writespace{1.3cm}{}

      \textbf{(b)} Give the absolute maximum of $W$ and when it occurs. Interpret it in context,
      with units.
      \writespace{1.6cm}{}

      \textbf{(c)} A dock worker says, ``The low tide is 9.'' Explain what the worker has
      confused, and state the low tide correctly in context.
      \writespace{1.8cm}{}

      \textbf{(d)} On the interval $(3, 6)$, is $W$ increasing or decreasing, and is it concave up
      or concave down? Describe in words what the water is doing.
      \writespace{1.8cm}{}

      \textbf{(e)} At what time is the water level falling fastest? Use concavity and the point
      of inflection to explain.
      \writespace{1.8cm}{}

\end{enumerate}

\end{document}
